\documentclass[11pt,a4paper]{article}
\usepackage[utf8]{inputenc}
\usepackage[T1]{fontenc}
\usepackage[spanish,es-minimal,es-noshorthands]{babel}
\usepackage{lmodern}
\linespread{1.22}
\usepackage[a4paper,top=2.4cm,bottom=2.2cm,left=1.5cm,right=1.5cm,headheight=1.1cm,headsep=0.5cm,footskip=1cm]{geometry}
\usepackage{amsmath,amssymb,cancel}
\usepackage{icomma}
\usepackage{graphicx}
\usepackage[table]{xcolor}
\usepackage{tikz}\usetikzlibrary{calc}
\usepackage{pgfplots}\pgfplotsset{compat=1.17}
\usepackage{fontawesome5}
\usepackage{ccicons}
\usepackage{fancyhdr}
\usepackage{enumitem}
\usepackage{booktabs}
\usepackage{tabularx}
\usepackage{needspace}
\usepackage[most]{tcolorbox}
\usepackage{titlesec}
\usepackage[hidelinks]{hyperref}
\hypersetup{pdftitle={Apuntes de teoría. Unidad 3: Sistemas de ecuaciones e inecuaciones. 1.º de Bachillerato CCSS},pdfsubject={Matemáticas Aplicadas a las Ciencias Sociales I},pdfkeywords={sistemas de ecuaciones, método de Gauss, sistemas no lineales, inecuaciones, sistemas de inecuaciones, inecuaciones con dos incógnitas, semiplanos}}

% colors de la web (FentMates, bloque de Álgebra: azul; es manté el nom fmgreen dels apunts de les altres unitats)
\definecolor{fmnavy}{HTML}{1B2748}
\definecolor{fmgreen}{HTML}{3D6AA0}
\definecolor{fmgreenl}{HTML}{4E81BD}
\definecolor{fmyellow}{HTML}{FFC94D}
\definecolor{fmgold}{HTML}{E9A919}
\definecolor{fmaccent}{HTML}{FF7A50}
\definecolor{fmteal}{HTML}{2BB3A3}
\definecolor{fmdef}{HTML}{1D4F73}\definecolor{fmdefbg}{HTML}{EEF4F9}
\definecolor{fmprop}{HTML}{3B4252}\definecolor{fmpropbg}{HTML}{F3F4F7}
\definecolor{fmex}{HTML}{A3203A}\definecolor{fmexbg}{HTML}{FBF0F2}
\definecolor{fmmuted}{HTML}{5D6685}\definecolor{fmball0}{HTML}{E2483D}\definecolor{fmball1}{HTML}{2F6FD6}
\definecolor{fmsurface}{HTML}{EEF1F8}
\definecolor{fmA}{HTML}{3B82F6}\definecolor{fmB}{HTML}{F97316}\definecolor{fmR}{HTML}{E11D48}\definecolor{fmK}{HTML}{16A34A}\definecolor{fmV}{HTML}{7C3AED}

% blocs que no es poden partir: es mesuren i es reserva exactament el seu espai (amb el títol o la capçalera del requadre)
\newsavebox{\fmkb}
\newcommand{\fmmeasure}[2]{\setbox\fmkb=\vbox{\hsize=\dimexpr\linewidth-#1\relax\linewidth=\hsize #2}}
\newcommand{\fmblock}[1]{\fmmeasure{0pt}{#1}\Needspace{\dimexpr\ht\fmkb+\dp\fmkb\relax}#1}
\newcommand{\fmheadblock}[2]{\fmmeasure{0pt}{#2}\Needspace{\dimexpr\ht\fmkb+\dp\fmkb+2.2cm\relax}#1#2}
\newcommand{\fmboxneed}[1]{\fmmeasure{9mm}{#1}\Needspace{\dimexpr\ht\fmkb+\dp\fmkb+1.6cm\relax}}
\newif\iffmstack\newif\iffmfit
\newcommand{\fmfits}[1]{\par\ifdim\dimexpr\pagegoal-\pagetotal\relax<\dimexpr\ht\fmkb+\dp\fmkb+#1\relax\global\fmfitfalse\else\global\fmfittrue\fi}
\newcommand{\fmtwo}[2]{\fmmeasure{0pt}{#1}\fmfits{0pt}\iffmfit #1\else #2\fi}
\newcommand{\fmheadtwo}[3]{\fmmeasure{0pt}{#2}\fmfits{2.2cm}\iffmfit #1#2\else\Needspace{\dimexpr1.6cm+7\baselineskip\relax}#1#3\fi}
\newcommand{\fmtwoin}[2]{\iffmstack #2\else #1\fi}
\newcommand{\fmboxneedalt}[1]{\fmmeasure{9mm}{#1}\fmfits{1.6cm}\iffmfit\global\fmstackfalse\else\global\fmstacktrue\Needspace{8\baselineskip}\fi}
\newcommand{\fmheadboxalt}[2]{\fmmeasure{9mm}{#2}\fmfits{3.4cm}\iffmfit\global\fmstackfalse\else\global\fmstacktrue\Needspace{\dimexpr2.4cm+5\baselineskip\relax}\fi#1}
\newcommand{\fmheadbox}[2]{\fmmeasure{9mm}{#2}\Needspace{\dimexpr\ht\fmkb+\dp\fmkb+3.4cm\relax}#1}
\newcommand{\lt}{<}\newcommand{\gt}{>}
% fórmula en línia que no es parteix mai; si no cap en una línia, es posa centrada
% fórmula en línia; si ocupa més del 80 % de la línia, va centrada (amb la puntuació que la segueix)
\makeatletter
\newcommand{\fmm}[1]{\setbox0=\hbox{$#1$}\ifdim\wd0>0.8\linewidth\def\fmmnext{\fmmdisp{#1}}\else\def\fmmnext{\mbox{\box0}}\fi\fmmnext}
\def\fmmdisp#1{\@ifnextchar.{\fmmdispp{#1}}{\@ifnextchar,{\fmmdispc{#1}}{\@ifnextchar;{\fmmdisps{#1}}{\fmdisp{#1}}}}}
\def\fmmdispp#1.{\fmdisp{#1\,.}}\def\fmmdispc#1,{\fmdisp{#1\,,}}\def\fmmdisps#1;{\fmdisp{#1\,;}}
\makeatother
% fórmula centrada; si és més ampla que la línia, s'encongeix fins que hi cap
\newcommand{\fmdisp}[1]{\[\setbox2=\hbox{$\displaystyle #1$}\ifdim\wd2>\linewidth\resizebox{\linewidth}{!}{\box2}\else\box2\fi\]}
\DeclareUnicodeCharacter{2212}{-}
\DeclareUnicodeCharacter{2225}{\ensuremath{\parallel}}\DeclareUnicodeCharacter{22A5}{\ensuremath{\perp}}\DeclareUnicodeCharacter{2226}{\ensuremath{\nparallel}}
\DeclareUnicodeCharacter{207A}{\ensuremath{^{+}}}\DeclareUnicodeCharacter{207B}{\ensuremath{^{-}}}
\setlength{\parindent}{0pt}
\setlength{\parskip}{12pt}
\setlength{\lineskip}{5pt}\setlength{\lineskiplimit}{5pt}  % més aire al voltant de les fraccions en línia
\setlength{\emergencystretch}{3em}
\setlist{itemsep=4pt,topsep=4pt}
\renewcommand{\thesection}{3.\arabic{section}}
\titleformat{\section}{\color{fmnavy}\Large\bfseries}{\colorbox{fmdef}{\color{white}\strut\ \thesection\ }}{0.6em}{}[{\color{fmdef}\titlerule[1.2pt]}]
\titlespacing*{\section}{0pt}{10pt}{4pt}
\titleformat*{\subsection}{\color{fmgreen}\large\bfseries}
\titlespacing*{\subsection}{0pt}{6pt plus 2pt}{0pt}
\allowdisplaybreaks

% dibuix de la targeta de la unitat a la web (illustrations.js), en blanc
\newcommand{\fmcard}{\begin{tikzpicture}[x=0.062cm,y=-0.062cm,white,line width=4pt,line cap=round]
  \begin{scope}[transparency group,opacity=0.15]
  \draw (20,106)--(220,20);
  \draw[opacity=0.75] (20,26)--(220,96);
  \draw[opacity=0.55] (128,6)--(128,114);
  \fill (128,59.5) circle (7);
  \node[anchor=base west,inner sep=0pt,font=\fontsize{21}{21}\selectfont] at (18,124) {$x+y+z=6$};
  \node[anchor=base west,inner sep=0pt,font=\fontsize{21}{21}\selectfont] at (150,12) {$2x-y=1$};
  \end{scope}
\end{tikzpicture}}
\newcommand{\fmbrand}{\raisebox{-0.25\height}{\includegraphics[height=0.5cm]{logo-fentmates.pdf}}}

% capçalera i peu
\pagestyle{fancy}\fancyhf{}
\fancyhead[L]{\small\fmbrand}
\fancyhead[R]{\small\color{fmmuted}1.º Bachillerato CCSS · Unidad 3: Sistemas de ecuaciones e inecuaciones}
\fancyfoot[L]{\scriptsize\color{fmmuted}\href{https://creativecommons.org/licenses/by-nc-sa/4.0/deed.ca}{\ccbyncsa}\ \ FentMates · fentmates.org · CC BY-NC-SA 4.0}
\fancyfoot[R]{\small\color{fmmuted}\thepage}
\renewcommand{\headrulewidth}{0.4pt}\renewcommand{\headrule}{\color{fmmuted!40}\hrule height \headrulewidth}
\fancypagestyle{primera}{\fancyhf{}\renewcommand{\headrulewidth}{0pt}
  \fancyfoot[L]{\scriptsize\color{fmmuted}\href{https://creativecommons.org/licenses/by-nc-sa/4.0/deed.ca}{\ccbyncsa}\ \ FentMates · fentmates.org · CC BY-NC-SA 4.0}
  \fancyfoot[R]{\small\color{fmmuted}\thepage}}

% caixes (mateixos colors que la web)
% quan una caixa es parteix entre dues pàgines, la vora del tall és discontínua
\tcbset{fmbase/.style={enhanced,breakable,arc=3mm,boxrule=1pt,left=3.5mm,right=3.5mm,top=2mm,bottom=2.5mm,before skip=10pt,after skip=10pt,fonttitle=\bfseries,
  frame style={fill=tcbcolframe},
  overlay first={\fill[tcbcolback] ([xshift=1pt]frame.south west) rectangle ([xshift=-1pt,yshift=1.2pt]frame.south east);
    \draw[tcbcolframe,line width=1pt,dash pattern=on 5pt off 3pt] (frame.south west)--(frame.south east);},
  overlay middle={\fill[tcbcolback] ([xshift=1pt]frame.south west) rectangle ([xshift=-1pt,yshift=1.2pt]frame.south east);
    \fill[tcbcolback] ([xshift=1pt]frame.north west) rectangle ([xshift=-1pt,yshift=-1.2pt]frame.north east);
    \draw[tcbcolframe,line width=1pt,dash pattern=on 5pt off 3pt] (frame.south west)--(frame.south east) (frame.north west)--(frame.north east);},
  overlay last={\fill[tcbcolback] ([xshift=1pt]frame.north west) rectangle ([xshift=-1pt,yshift=-1.2pt]frame.north east);
    \draw[tcbcolframe,line width=1pt,dash pattern=on 5pt off 3pt] (frame.north west)--(frame.north east);}}}
\newtcolorbox{fmdef}[1]{fmbase,colback=fmdefbg,colframe=fmdef,coltitle=white,title={#1}}
\newtcolorbox{fmprop}[1]{fmbase,colback=fmpropbg,colframe=fmprop,coltitle=white,title={#1},halign title=center}
\newtcolorbox{fmex}[1]{fmbase,before upper=\raggedright,colback=fmexbg,colframe=fmex,colbacktitle=fmexbg,coltitle=fmex,title={\faPencil*\ #1},
  titlerule=0.4pt,titlerule style={fmex!40,dashed}}
\newtcolorbox{fmwarn}[1]{fmbase,colback=fmsf!6,colframe=fmsf,coltitle=white,title={#1}}
\newtcolorbox{fmnote}[1][]{enhanced,breakable,colback=white,colframe=white,borderline west={3pt}{0pt}{fmaccent},arc=0pt,left=3mm,right=2mm,top=1mm,bottom=1mm,before skip=6pt,after skip=6pt,before upper=\raggedright,#1}
\newcommand{\fmkey}[1]{\colorbox{fmyellow!45}{\strut #1}}
\newcommand{\fmuline}[1]{\textbf{\color{fmaccent}#1}}
\definecolor{fmok}{HTML}{15803D}\definecolor{fmev}{HTML}{2563EB}\definecolor{fmsf}{HTML}{C2410C}\definecolor{fmsi}{HTML}{7C3AED}\definecolor{fmas}{HTML}{A3203A}
\newcommand{\fmbadge}[2]{\tcbox[on line,colback=#1,colframe=#1,arc=2mm,boxsep=0pt,left=4pt,right=4pt,top=1.5pt,bottom=1.5pt]{\color{white}\footnotesize\bfseries #2}}
\newcommand{\fmchip}[1]{\tcbox[on line,colback=fmsf!10,colframe=fmsf!50,arc=2mm,boxsep=1pt,left=4pt,right=4pt]{#1}}
\newcommand{\fmlevel}[2]{\tcbox[on line,colback=white,colframe=fmmuted!40,arc=2mm,boxsep=2pt,halign=center]{\begin{tabular}{c}{\scriptsize\color{fmmuted}\MakeUppercase{#1}}\\#2\end{tabular}}}
% gràfiques
\pgfplotsset{fmaxis/.style={axis lines=middle,grid=major,grid style={fmmuted!15},axis line style={fmmuted,-latex},
  tick label style={font=\scriptsize,color=fmmuted},clip=true,unbounded coords=jump,restrict y to domain=-200:200,
  every axis plot/.append style={line join=round}},
  fmcurve/.style={fmaccent,line width=1.6pt,smooth}}
\tikzset{fmasy/.style={fmmuted,dashed,line width=0.9pt},fmguide/.style={fmaccent,dotted,line width=0.8pt},
  fmdot/.style={circle,fill=fmaccent,inner sep=0pt,minimum size=5pt},fmopen/.style={circle,draw=fmaccent,fill=white,line width=1pt,inner sep=0pt,minimum size=5pt}}
\newcommand{\fmpau}[1]{\tcbox[on line,colback=fmnavy,colframe=fmnavy,arc=2mm,boxsep=0pt,left=4pt,right=4pt,top=1.5pt,bottom=1.5pt]{\color{white}\scriptsize\bfseries\MakeUppercase{#1}}\ }

\begin{document}
\thispagestyle{primera}
\begin{tikzpicture}[remember picture,overlay]
  \fill[fmgreen] (current page.north west) rectangle ([yshift=-10.1cm]current page.north east);
  \begin{scope}
    \clip (current page.north west) rectangle ([yshift=-10.1cm]current page.north east);
    \node[anchor=north west,inner sep=0pt] at ([xshift=6.6cm,yshift=-1.3cm]current page.north west) {\fmcard};
  \end{scope}
  \node[anchor=north west,inner sep=0pt] at ([xshift=1.5cm,yshift=-1.1cm]current page.north west) {\tcbox[colback=white,colframe=white,arc=2.5mm,boxsep=0pt,left=7pt,right=7pt,top=5pt,bottom=5pt]{\includegraphics[height=0.75cm]{logo-fentmates.pdf}}};
\end{tikzpicture}
\vspace*{0.0cm}
\tcbox[colback=white!15!fmgreen,colframe=white!15!fmgreen,arc=2.5mm,boxsep=0pt,left=6pt,right=6pt,top=3pt,bottom=3pt]{\color{white}\small\bfseries Apuntes de teoría · 1.º Bachillerato · Matemáticas Aplicadas a las Ciencias Sociales I}

\vspace{4pt}
{\raggedright\color{white}\fontsize{25}{30}\selectfont\bfseries{\color{fmyellow}Unidad 3} · Sistemas de ecuaciones e inecuaciones\par}
\vspace{6pt}
\begin{minipage}{\linewidth}\color{white!88!fmgreen}\small Teoría y ejemplos resueltos paso a paso: sistemas de ecuaciones lineales con paréntesis y denominadores, el método de Gauss, sistemas no lineales, inecuaciones de primer grado, de segundo grado y de grado superior, sistemas de inecuaciones con una incógnita, inecuaciones y sistemas con dos incógnitas y problemas de programación lineal.\end{minipage}\par
% l'índex comença sempre a la mateixa altura, per davall de la franja (com en les unitats 5-10)
\vspace*{\dimexpr 8.6cm-\pagetotal\relax}

\begin{tcolorbox}[colback=fmsurface,colframe=fmsurface,arc=3mm,left=4mm,right=4mm,top=2mm,bottom=2mm]
\small\textbf{Índice}\par\smallskip
\begin{minipage}[t]{0.49\linewidth}\small \mbox{\hyperlink{sec1}{\textbf{\color{fmgreen}3.1}\ \ Sistemas de ecuaciones lineales}}\\[3pt]
\mbox{\hyperlink{sec2}{\textbf{\color{fmgreen}3.2}\ \ Método de Gauss}}\\[3pt]
\mbox{\hyperlink{sec3}{\textbf{\color{fmgreen}3.3}\ \ Sistemas de ecuaciones no lineales}}\\[3pt]
\mbox{\hyperlink{sec4}{\textbf{\color{fmgreen}3.4}\ \ Inecuaciones de primer grado}}\\[3pt]
\hyperlink{sec5}{\textbf{\color{fmgreen}3.5}\ \ Inecuaciones de segundo grado y de grado superior}\\[3pt]\end{minipage}\hfill\begin{minipage}[t]{0.49\linewidth}\small \hyperlink{sec6}{\textbf{\color{fmgreen}3.6}\ \ Sistemas de inecuaciones con una incógnita}\\[3pt]
\mbox{\hyperlink{sec7}{\textbf{\color{fmgreen}3.7}\ \ Inecuaciones con dos incógnitas}}\\[3pt]
\hyperlink{sec8}{\textbf{\color{fmgreen}3.8}\ \ Sistemas de inecuaciones con dos incógnitas}\\[3pt]
\mbox{\hyperlink{sec9}{\textbf{\color{fmgreen}3.9}\ \ Programación lineal}}\\[3pt]\end{minipage}
\end{tcolorbox}

\Needspace{5.7cm}
\section{Sistemas de ecuaciones lineales}\hypertarget{sec1}{}
\begin{fmdef}{Sistemas y soluciones}
\begin{itemize}
\item Un \textbf{sistema de ecuaciones lineales} es un conjunto de ecuaciones de primer grado con varias incógnitas. Una \textbf{solución} es un conjunto de valores de las incógnitas que verifica \textbf{todas} las ecuaciones a la vez.
\item Con dos incógnitas, cada ecuación \fmm{ax+by=c} es una \textbf{recta} del plano, y las soluciones del sistema son los puntos comunes a las dos rectas.
\item Dos sistemas son \textbf{equivalentes} si tienen las mismas soluciones.
\end{itemize}

\begin{center}\renewcommand{\arraystretch}{1.5}\setlength{\parskip}{0pt}\small\begin{tabularx}{\linewidth}{>{\raggedright\arraybackslash}p{4.2cm}>{\raggedright\arraybackslash}X>{\raggedright\arraybackslash}X>{\raggedright\arraybackslash}X}\toprule
\textbf{Tipo} & \textbf{Soluciones} & \textbf{Las rectas\dots{}} & \textbf{Coeficientes}\\\midrule
\fmbadge{fmok}{SCD} Compatible determinado & una única solución & se cortan en un punto & $\dfrac a{a'}\neq\dfrac b{b'}$\\[3pt]\end{tabularx}\par\nointerlineskip\vskip0pt
\small\begin{tabularx}{\linewidth}{>{\raggedright\arraybackslash}p{4.2cm}>{\raggedright\arraybackslash}X>{\raggedright\arraybackslash}X>{\raggedright\arraybackslash}X}\fmbadge{fmev}{SCI} Compatible indeterminado & infinitas soluciones & son coincidentes & $\dfrac a{a'}=\dfrac b{b'}=\dfrac c{c'}$\\[3pt]\end{tabularx}\par\nointerlineskip\vskip0pt
\small\begin{tabularx}{\linewidth}{>{\raggedright\arraybackslash}p{4.2cm}>{\raggedright\arraybackslash}X>{\raggedright\arraybackslash}X>{\raggedright\arraybackslash}X}\fmbadge{fmas}{SI} Incompatible & ninguna solución & son paralelas & $\dfrac a{a'}=\dfrac b{b'}\neq\dfrac c{c'}$\\[3pt]\bottomrule\end{tabularx}\end{center}
\end{fmdef}

\Needspace{5\baselineskip}
\begin{fmprop}{Métodos de resolución}
\begin{itemize}
\item \textbf{Sustitución}: se despeja una incógnita en una ecuación y se sustituye en la otra.
\item \textbf{Igualación}: se despeja la misma incógnita en las dos ecuaciones y se igualan las expresiones.
\item \textbf{Reducción}: se multiplican las ecuaciones por números adecuados para que, al sumarlas o restarlas, desaparezca una incógnita.
\end{itemize}

Cuando las ecuaciones tienen \textbf{paréntesis} o \textbf{denominadores}, primero las simplificamos hasta la forma \fmm{ax+by=c}, como en las ecuaciones de primer grado:

\begin{enumerate}
\item Quitamos los paréntesis (atención a los signos \fmm{-} delante de un paréntesis).
\item Quitamos los denominadores multiplicando los dos miembros por el m.c.m. de los denominadores.
\item Pasamos las incógnitas a la izquierda y los números a la derecha, y reducimos.
\end{enumerate}
\end{fmprop}

\Needspace{5\baselineskip}
\begin{fmex}{Ejemplo resuelto: sistemas con paréntesis y denominadores}
\begin{enumerate}[label=\textbf{\arabic*.},leftmargin=*]
\item Resuelve \fmm{\begin{cases}3(x-1)-2(y+1)=1\\ 2(x+y)-(x-3)=13\end{cases}}

Quitamos paréntesis: \fmm{3x-3-2y-2=1\Rightarrow 3x-2y=6};\fmm{\quad 2x+2y-x+3=13\Rightarrow x+2y=10}

Reducción: sumamos las dos ecuaciones y desaparece la \fmm{y}: \fmm{4x=16\Rightarrow x=4}

Sustituimos en la segunda: \fmm{4+2y=10\Rightarrow y=3}. Solución: \fmm{(x,y)=(4,3)} \fmbadge{fmok}{SCD}
\item Resuelve \fmm{\begin{cases}\dfrac x2+\dfrac y3=4\\[2mm] \dfrac{x+2}3-\dfrac{y-2}4=1\end{cases}}

Primera ecuación por 6 (m.c.m. de 2 y 3): \fmm{3x+2y=24}

Segunda ecuación por 12: \fmm{4(x+2)-3(y-2)=12\Rightarrow 4x+8-3y+6=12\Rightarrow 4x-3y=-2}

Reducción: la primera por 3 y la segunda por 2, y sumamos: \fmm{9x+6y=72,\ \ 8x-6y=-4\Rightarrow 17x=68\Rightarrow x=4}

\fmm{3\cdot4+2y=24\Rightarrow y=6}. Solución: \fmm{(x,y)=(4,6)}
\item Clasifica \fmm{\begin{cases}2x-y=3\\ 4x-2y=6\end{cases}} y \fmm{\begin{cases}2x-y=3\\ 4x-2y=1\end{cases}}

En el primero, la segunda ecuación es la primera multiplicada por 2: \fmm{\frac24=\frac{-1}{-2}=\frac36}. Son la misma recta: \textbf{SCI}. Si hacemos \fmm{x=\lambda}, \fmm{y=2\lambda-3}: \fmm{(x,y)=(\lambda,\ 2\lambda-3)}.

En el segundo, \fmm{\frac24=\frac{-1}{-2}\neq\frac31}: las rectas son paralelas. Por reducción queda \fmm{0=5}: \textbf{SI}, no tiene solución.
\end{enumerate}
\end{fmex}

\Needspace{5.7cm}
\section{Método de Gauss}\hypertarget{sec2}{}
\begin{fmdef}{Sistemas de tres ecuaciones con tres incógnitas}
Con tres incógnitas \fmm{x,y,z} los sistemas se clasifican igual que antes (\textbf{SCD}, \textbf{SCI} o \textbf{SI}), pero resolverlos por sustitución es largo. El método de Gauss lo hace de manera ordenada.

Un sistema es \textbf{escalonado} (o triangular) si cada ecuación tiene una incógnita menos que la anterior. Se resuelve de abajo hacia arriba:

\fmdisp{\begin{cases}x+y+z=6\\ \quad\ \ y+z=5\\ \qquad\ \ \, 2z=6\end{cases}\Rightarrow z=3,\ \ y=5-3=2,\ \ x=6-2-3=1}

Obtenemos un sistema \textbf{equivalente} (con las mismas soluciones) si intercambiamos dos ecuaciones, multiplicamos una ecuación por un número \fmm{\neq0} o sumamos a una ecuación un múltiplo de otra.
\end{fmdef}

\Needspace{5\baselineskip}
\begin{fmprop}{Método de Gauss}
Consiste en transformar el sistema en otro equivalente \textbf{escalonado}. Trabajamos con la \textbf{matriz ampliada} (coeficientes y términos independientes):

\fmdisp{\left(\begin{array}{ccc|c}a&b&c&d\\ a'&b'&c'&d'\\ a''&b''&c''&d''\end{array}\right)\longrightarrow\left(\begin{array}{ccc|c}*&*&*&*\\ 0&*&*&*\\ 0&0&m&n\end{array}\right)}

\begin{enumerate}
\item Hacemos ceros en la primera columna, debajo del primer elemento, combinando cada fila con la primera: \fmm{F_2\to aF_2-a'F_1}\dots{}
\item Hacemos ceros en la segunda columna debajo del segundo elemento, combinando con la segunda fila.
\item Miramos la última fila \fmm{\left(0\ \ 0\ \ m\,|\,n\right)}:

\begin{itemize}
\item \fmm{m\neq0}: \textbf{SCD}. Se despeja \fmm{z=\frac nm} y se sustituye hacia arriba.
\item \fmm{m=0,\ n=0}: \textbf{SCI}. La última ecuación desaparece (\fmm{0=0}) y una incógnita pasa a ser un parámetro (\fmm{z=\lambda}).
\item \fmm{m=0,\ n\neq0}: \textbf{SI}. La última ecuación es \fmm{0=n}, imposible.
\end{itemize}
\end{enumerate}

\begin{fmnote}\fmuline{Consejo}: conviene poner en primer lugar una ecuación con coeficiente 1 (o \fmm{-1}) en la primera incógnita.\end{fmnote}
\end{fmprop}

\Needspace{5\baselineskip}
\begin{fmex}{Ejemplo resuelto: método de Gauss}
\begin{enumerate}[label=\textbf{\arabic*.},leftmargin=*]
\item Resuelve \fmm{\begin{cases}x+y+z=6\\ 2x-y+z=3\\ x+2y-z=2\end{cases}}

\fmm{\left(\begin{array}{rrr|r}1&1&1&6\\ 2&-1&1&3\\ 1&2&-1&2\end{array}\right)\xrightarrow[F_3\to F_3-F_1]{F_2\to F_2-2F_1}\left(\begin{array}{rrr|r}1&1&1&6\\ 0&-3&-1&-9\\ 0&1&-2&-4\end{array}\right)}

\fmm{\xrightarrow{F_3\to 3F_3+F_2}\left(\begin{array}{rrr|r}1&1&1&6\\ 0&-3&-1&-9\\ 0&0&-7&-21\end{array}\right)}

\fmm{-7z=-21\Rightarrow z=3;\quad -3y-3=-9\Rightarrow y=2;\quad x+2+3=6\Rightarrow x=1}. \textbf{SCD}: \fmm{(x,y,z)=(1,2,3)}
\item Resuelve \fmm{\begin{cases}x-y+z=1\\ 2x+y-z=5\\ x+2y-2z=4\end{cases}}

\fmm{\left(\begin{array}{rrr|r}1&-1&1&1\\ 2&1&-1&5\\ 1&2&-2&4\end{array}\right)\xrightarrow[F_3\to F_3-F_1]{F_2\to F_2-2F_1}\left(\begin{array}{rrr|r}1&-1&1&1\\ 0&3&-3&3\\ 0&3&-3&3\end{array}\right)\xrightarrow{F_3\to F_3-F_2}\left(\begin{array}{rrr|r}1&-1&1&1\\ 0&3&-3&3\\ 0&0&0&0\end{array}\right)}

La última fila es \fmm{0=0}: \textbf{SCI}. Hacemos \fmm{z=\lambda}: \fmm{\ 3y-3\lambda=3\Rightarrow y=1+\lambda;\ \ x=1+y-z=2}

\fmm{(x,y,z)=(2,\ 1+\lambda,\ \lambda),\ \lambda\in\mathbb R}. Si el término independiente de la tercera ecuación fuera 5, la última fila quedaría \fmm{\left(0\ \ 0\ \ 0\,|\,1\right)}: \fmm{0=1}, el sistema sería \textbf{incompatible}.
\item Un grupo compra 10 entradas para un museo: las de adulto cuestan 8 €, las de joven 6 € y las infantiles 4 €. Pagan 56 € y hay tantas entradas infantiles como de adulto y de joven juntas. ¿Cuántas han comprado de cada tipo?

Sean \fmm{x}, \fmm{y}, \fmm{z} las entradas de adulto, de joven e infantiles: \fmm{\begin{cases}x+y+z=10\\ 8x+6y+4z=56\\ x+y-z=0\end{cases}}

\fmm{\left(\begin{array}{rrr|r}1&1&1&10\\ 8&6&4&56\\ 1&1&-1&0\end{array}\right)\xrightarrow[F_3\to F_3-F_1]{F_2\to F_2-8F_1}\left(\begin{array}{rrr|r}1&1&1&10\\ 0&-2&-4&-24\\ 0&0&-2&-10\end{array}\right)}

\fmm{z=5;\quad -2y-20=-24\Rightarrow y=2;\quad x=10-2-5=3}. 3 de adulto, 2 de joven y 5 infantiles.
\end{enumerate}
\end{fmex}

\Needspace{5.7cm}
\section{Sistemas de ecuaciones no lineales}\hypertarget{sec3}{}
\begin{fmdef}{Sistemas no lineales}
Son sistemas en los que alguna ecuación no es de primer grado: aparecen cuadrados, productos de incógnitas (\fmm{xy}), incógnitas en el denominador\dots{} Los métodos más habituales son:

\begin{itemize}
\item \textbf{Sustitución}: se despeja una incógnita (preferiblemente de una ecuación lineal) y se sustituye en la otra. Suele quedar una ecuación de segundo grado.
\item \textbf{Reducción}: si las incógnitas aparecen con la misma forma (\fmm{x^2} e \fmm{y^2}, por ejemplo), se combinan las ecuaciones para eliminar una.
\item Si hay \textbf{denominadores} con incógnitas, se multiplica por el m.c.m.; si hay \textbf{radicales}, se aísla el radical y se eleva al cuadrado. En ambos casos pueden aparecer soluciones que no lo son.
\end{itemize}

Cada solución es una \textbf{pareja} \fmm{(x,y)}: por cada valor de una incógnita hay que calcular el valor correspondiente de la otra. Gráficamente, las soluciones son los puntos de corte de las dos curvas.

\begin{fmnote}Conviene \textbf{comprobar} las soluciones en las dos ecuaciones, sobre todo si hay denominadores o raíces.\end{fmnote}
\end{fmdef}

\Needspace{5\baselineskip}
\begin{fmex}{Ejemplo resuelto: sistemas no lineales}
\begin{enumerate}[label=\textbf{\arabic*.},leftmargin=*]
\item \fmm{\begin{cases}x-y=1\\ x^2+y^2=13\end{cases}}

Sustitución: \fmm{x=1+y\Rightarrow(1+y)^2+y^2=13\Rightarrow 2y^2+2y-12=0\Rightarrow y^2+y-6=0}

\fmm{y=2\Rightarrow x=3;\quad y=-3\Rightarrow x=-2}. Soluciones: \fmm{(3,2)} y \fmm{(-2,-3)}
\item \fmm{\begin{cases}x+y=5\\ xy=6\end{cases}}

\fmm{y=5-x\Rightarrow x(5-x)=6\Rightarrow x^2-5x+6=0\Rightarrow x=2} o \fmm{x=3}

\fmm{x=2\Rightarrow y=3;\quad x=3\Rightarrow y=2}. Soluciones: \fmm{(2,3)} y \fmm{(3,2)}
\item \fmm{\begin{cases}x^2+y^2=25\\ x^2-y^2=7\end{cases}}

Reducción: sumamos las ecuaciones, \fmm{2x^2=32\Rightarrow x^2=16\Rightarrow x=\pm4}

\fmm{y^2=25-16=9\Rightarrow y=\pm3}. Hay cuatro soluciones: \fmm{(4,3),\ (4,-3),\ (-4,3),\ (-4,-3)}
\item Con denominadores: \fmm{\begin{cases}x+y=5\\[1mm] \dfrac1x+\dfrac1y=\dfrac56\end{cases}}

Sumamos las fracciones: \fmm{\dfrac{x+y}{xy}=\dfrac56\Rightarrow\dfrac{5}{xy}=\dfrac56\Rightarrow xy=6}

\fmm{y=5-x\Rightarrow x(5-x)=6\Rightarrow x^2-5x+6=0\Rightarrow x=2} o \fmm{x=3}. Soluciones: \fmm{(2,3)} y \fmm{(3,2)} (ninguna anula un denominador)
\item Con radicales: \fmm{\begin{cases}\sqrt x+y=5\\ x-y=7\end{cases}}

\fmm{x=y+7\Rightarrow\sqrt{y+7}=5-y}. Elevamos al cuadrado: \fmm{y+7=25-10y+y^2\Rightarrow y^2-11y+18=0\Rightarrow y=2} o \fmm{y=9}

Comprobamos: con \fmm{y=2}, \fmm{x=9} y \fmm{\sqrt9+2=5} . Con \fmm{y=9}, \fmm{x=16} y \fmm{\sqrt{16}+9=13\neq5}  (al elevar al cuadrado ha aparecido una solución falsa). Solución: \fmm{(9,2)}
\item El perímetro de un rectángulo mide 14 cm y su área, 12 cm$^2$. ¿Cuánto miden los lados?

\fmm{\begin{cases}2x+2y=14\\ xy=12\end{cases}\Rightarrow y=7-x\Rightarrow x(7-x)=12\Rightarrow x^2-7x+12=0\Rightarrow x=3} o \fmm{x=4}

Los lados miden 3 cm y 4 cm.
\end{enumerate}
\end{fmex}

\Needspace{5.7cm}
\section{Inecuaciones de primer grado}\hypertarget{sec4}{}
\begin{fmdef}{Inecuaciones y soluciones}
\begin{itemize}
\item Una \textbf{inecuación} es una desigualdad entre dos expresiones algebraicas, con uno de los signos \fmm{\lt,\ \gt,\ \le,\ \ge}. Por ejemplo, \fmm{3x-5\le x+7}.
\item Una \textbf{solución} es un valor de \fmm{x} que hace cierta la desigualdad. Normalmente hay infinitas y forman un \textbf{intervalo} o una unión de intervalos. \textbf{Resolver} una inecuación es encontrarlas todas.
\item Una inecuación es de \textbf{primer grado} si, después de simplificar, la \fmm{x} solo aparece elevada a 1: se puede escribir como \fmm{ax\lt b} (o con \fmm{\gt,\le,\ge}).
\end{itemize}

\begin{center}\renewcommand{\arraystretch}{1.5}\small\begin{tabularx}{\linewidth}{>{\raggedright\arraybackslash}p{3.2cm}>{\raggedright\arraybackslash}X>{\raggedright\arraybackslash}X>{\raggedright\arraybackslash}X>{\raggedright\arraybackslash}X>{\raggedright\arraybackslash}X}\toprule
\textbf{Desigualdad} & \textbf{$x\lt a$} & \textbf{$x\le a$} & \textbf{$x\gt a$} & \textbf{$x\ge a$} & \textbf{$a\lt x\le b$}\\\midrule
Intervalo & $(-\infty,a)$ & $(-\infty,a]$ & $(a,+\infty)$ & $[a,+\infty)$ & $(a,b]$\\[3pt]
\bottomrule\end{tabularx}\end{center}

\begin{fmnote}El corchete \fmm{[\ ]} indica que el extremo \textbf{se incluye} (punto lleno en la recta) y el paréntesis \fmm{(\ )}, que \textbf{no se incluye} (punto vacío). En \fmm{\pm\infty} siempre va paréntesis.\end{fmnote}
\end{fmdef}

\Needspace{5\baselineskip}
\begin{fmprop}{Reglas para transformar una inecuación}
\begin{itemize}
\item Si \textbf{sumamos o restamos} el mismo número (o la misma expresión) a los dos miembros, la desigualdad \textbf{no cambia}: \fmm{x+3\lt5\iff x\lt2}.
\item Si \textbf{multiplicamos o dividimos} los dos miembros por un número \textbf{positivo}, la desigualdad \textbf{no cambia}: \fmm{2x\lt6\iff x\lt3}.
\item Si \textbf{multiplicamos o dividimos} por un número \textbf{negativo}, la desigualdad \textbf{cambia de sentido}: \fmm{-2x\lt6\iff x\gt-3}. (Pensad que \fmm{2\lt3} pero \fmm{-2\gt-3}.)
\end{itemize}
\end{fmprop}

\Needspace{5\baselineskip}
\begin{fmprop}{Cómo se resuelve una inecuación de primer grado}
\begin{enumerate}
\item Quitamos los \textbf{paréntesis} y, si hay, los \textbf{denominadores} (multiplicando por el m.c.m., que es positivo).
\item \textbf{Transponemos}: los términos con \fmm{x} a un miembro y los números al otro, y \textbf{reducimos}.
\item \textbf{Despejamos} la \fmm{x} dividiendo por el coeficiente. \textbf{Si es negativo, la desigualdad cambia de sentido.}
\item Escribimos la solución como \textbf{intervalo} y la representamos en la recta real.
\end{enumerate}
\end{fmprop}

\Needspace{5\baselineskip}
\begin{fmex}{Ejemplo resuelto: inecuaciones de primer grado}
\begin{enumerate}[label=\textbf{\arabic*.},leftmargin=*]
\item Resuelve \fmm{3x-5\le x+7}.

\fmm{3x-x\le7+5\ \Rightarrow\ 2x\le12\ \Rightarrow\ x\le6}

Solución: \fmm{x\in(-\infty,\,6]}. El 6 se incluye porque la desigualdad es \fmm{\le}.
\item Resuelve \fmm{2(x-1)\gt5x+4}.

Paréntesis: \fmm{2x-2\gt5x+4\ \Rightarrow\ 2x-5x\gt4+2\ \Rightarrow\ -3x\gt6}

Dividimos entre \fmm{-3}, \textbf{negativo}: la desigualdad cambia de sentido. \fmm{\ x\lt\dfrac{6}{-3}=-2}

Solución: \fmm{x\in(-\infty,\,-2)}
\item Resuelve \fmm{\dfrac{x+1}{2}-\dfrac{x}{3}\ge1}.

Multiplicamos por \fmm{\text{m.c.m.}(2,3)=6}, que es positivo: \fmm{\ 3(x+1)-2x\ge6}

\fmm{3x+3-2x\ge6\ \Rightarrow\ x\ge3}. Solución: \fmm{x\in[3,+\infty)}
\item Resuelve \fmm{2(x+1)\ge2x-3}.

\fmm{2x+2\ge2x-3\ \Rightarrow\ 0x\ge-5\ \Rightarrow\ 0\ge-5}

La desigualdad \fmm{0\ge-5} es cierta siempre, sea cual sea \fmm{x}: la solución es \fmm{\mathbb R}. (Si hubiera quedado una desigualdad falsa, como \fmm{0\ge5}, no tendría solución.)
\end{enumerate}
\end{fmex}

\Needspace{5\baselineskip}
\begin{fmex}{Ejemplo resuelto: un problema}
\begin{enumerate}[label=\textbf{\arabic*.},leftmargin=*]
\item Una empresa de alquiler de coches cobra 40 € fijos más 0,25 € por kilómetro, y otra, 60 € fijos más 0,15 € por kilómetro. ¿A partir de cuántos kilómetros es más barata la segunda?

\fmm{x=} kilómetros recorridos. La segunda es más barata cuando \fmm{\ 60+0{,}15x\lt40+0{,}25x}

\fmm{0{,}15x-0{,}25x\lt40-60\ \Rightarrow\ -0{,}1x\lt-20\ \Rightarrow\ x\gt\dfrac{-20}{-0{,}1}=200}

La segunda empresa es más barata si se hacen \textbf{más de 200 km}.
\end{enumerate}
\end{fmex}

\Needspace{5.7cm}
\section{Inecuaciones de segundo grado y de grado superior}\hypertarget{sec5}{}
\begin{fmprop}{Método general: estudio del signo}
\begin{enumerate}
\item Pasamos todos los términos a un miembro para que quede \fmm{P(x)\gt0} (o \fmm{\lt0,\ \le0,\ \ge0}).
\item Calculamos las \textbf{raíces} de \fmm{P(x)}: en segundo grado, con la fórmula; en grado superior, sacando factor común o con Ruffini.
\item Las raíces dividen la recta real en intervalos. En cada intervalo \fmm{P(x)} tiene siempre el mismo signo: lo averiguamos \textbf{probando un valor} del intervalo o con una \textbf{tabla de signos} de los factores.
\item Elegimos los intervalos donde el signo es el que pide la inecuación. Las raíces se incluyen si la desigualdad es \fmm{\le} o \fmm{\ge}.
\end{enumerate}
\end{fmprop}

\Needspace{5\baselineskip}
\begin{fmprop}{Interpretación gráfica en segundo grado}
Resolver \fmm{ax^2+bx+c\gt0} es buscar dónde la parábola \fmm{y=ax^2+bx+c} está \textbf{por encima} del eje \fmm{X} (y \fmm{\lt0}, por debajo). Si \fmm{a\gt0} y hay dos raíces \fmm{x_1\lt x_2}:

\begin{itemize}
\item \fmm{ax^2+bx+c\lt0} \textbf{entre} las raíces: \fmm{(x_1,x_2)}.
\item \fmm{ax^2+bx+c\gt0} \textbf{fuera} de las raíces: \fmm{(-\infty,x_1)\cup(x_2,+\infty)}.
\item Si no hay raíces reales (\fmm{\Delta\lt0}), la parábola no corta el eje y el polinomio tiene \textbf{siempre el mismo signo}: la solución es \fmm{\mathbb R} o \fmm{\varnothing}.
\item Si \fmm{a\lt0}, conviene multiplicar por \fmm{-1} (cambiando el sentido de la desigualdad) para tener \fmm{a\gt0}.
\end{itemize}
\end{fmprop}

\Needspace{5\baselineskip}
\begin{fmex}{Ejemplo resuelto: inecuaciones de segundo grado}
\begin{enumerate}[label=\textbf{\arabic*.},leftmargin=*]
\item Resuelve \fmm{x^2-5x+6\le0}.

Raíces: \fmm{x=\dfrac{5\pm\sqrt{25-24}}{2}=\dfrac{5\pm1}{2}\Rightarrow x=2,\ x=3}

Probamos un valor de cada intervalo: \fmm{P(0)=6\gt0;\ \ P(2{,}5)=-0{,}25\lt0;\ \ P(4)=2\gt0}

Queremos \fmm{P(x)\le0}: el intervalo del medio, con las raíces incluidas. \fmm{\ x\in[2,\,3]}
\item Resuelve \fmm{-x^2+2x+3\lt0}.

Multiplicamos por \fmm{-1} y cambiamos el sentido: \fmm{\ x^2-2x-3\gt0}. Raíces: \fmm{-1} y \fmm{3}.

La parábola \fmm{y=x^2-2x-3} se abre hacia arriba: es positiva fuera de las raíces. \fmm{\ x\in(-\infty,-1)\cup(3,+\infty)}
\item Resuelve \fmm{x^2+x+3\gt0} y \fmm{x^2-6x+9\gt0}.

\fmm{x^2+x+3}: \fmm{\Delta=1-12=-11\lt0}. No tiene raíces y \fmm{P(0)=3\gt0}: es siempre positivo. \fmm{\ x\in\mathbb R}

\fmm{x^2-6x+9=(x-3)^2}, que es positivo siempre excepto en \fmm{x=3}, donde vale 0. \fmm{\ x\in(-\infty,3)\cup(3,+\infty)}
\end{enumerate}
\end{fmex}

\Needspace{5\baselineskip}
\begin{fmex}{Ejemplo resuelto: grado superior}
\begin{enumerate}[label=\textbf{\arabic*.},leftmargin=*]
\item Resuelve \fmm{x^3-x^2-6x\ge0}.

Factor común y segundo grado: \fmm{\ x(x^2-x-6)=x(x+2)(x-3)}. Raíces: \fmm{-2,\ 0,\ 3}.

\begin{center}\renewcommand{\arraystretch}{1.5}\begin{tabular}{cccccccc}\toprule
\textbf{} & \textbf{$(-\infty,-2)$} & \textbf{$-2$} & \textbf{$(-2,0)$} & \textbf{$0$} & \textbf{$(0,3)$} & \textbf{$3$} & \textbf{$(3,+\infty)$}\\\midrule
$x+2$ & \textcolor{fmas}{$\boldsymbol{-}$} & 0 & \textcolor{fmok}{\textbf{+}} &  & \textcolor{fmok}{\textbf{+}} &  & \textcolor{fmok}{\textbf{+}}\\
$x$ & \textcolor{fmas}{$\boldsymbol{-}$} &  & \textcolor{fmas}{$\boldsymbol{-}$} & 0 & \textcolor{fmok}{\textbf{+}} &  & \textcolor{fmok}{\textbf{+}}\\
$x-3$ & \textcolor{fmas}{$\boldsymbol{-}$} &  & \textcolor{fmas}{$\boldsymbol{-}$} &  & \textcolor{fmas}{$\boldsymbol{-}$} & 0 & \textcolor{fmok}{\textbf{+}}\\
$P(x)$ & \textcolor{fmas}{$\boldsymbol{-}$} & 0 & \textcolor{fmok}{\textbf{+}} & 0 & \textcolor{fmas}{$\boldsymbol{-}$} & 0 & \textcolor{fmok}{\textbf{+}}\\
\bottomrule\end{tabular}\end{center}

Queremos \fmm{P(x)\ge0}: \fmm{\ x\in[-2,\,0]\cup[3,+\infty)}
\item Resuelve \fmm{x^3-2x^2-x+2\gt0}.

Con Ruffini (probando los divisores de 2) encontramos las raíces \fmm{1,\ -1} y \fmm{2}: \fmm{\ P(x)=(x+1)(x-1)(x-2)}

Signos: \fmm{(-\infty,-1):-;\ \ (-1,1):+;\ \ (1,2):-;\ \ (2,+\infty):+}

\fmm{x\in(-1,\,1)\cup(2,+\infty)}
\item El beneficio de una empresa, en miles de euros, cuando fabrica \fmm{x} centenares de unidades es \fmm{B(x)=-x^2+12x-20}. ¿Entre qué cantidades tiene beneficios?

\fmm{-x^2+12x-20\gt0\iff x^2-12x+20\lt0}. Raíces: \fmm{x=2} y \fmm{x=10}.

Negativo entre las raíces: \fmm{x\in(2,\,10)}. Tiene beneficios si fabrica entre 200 y 1 000 unidades.
\end{enumerate}
\end{fmex}

\Needspace{5.7cm}
\section{Sistemas de inecuaciones con una incógnita}\hypertarget{sec6}{}
\begin{fmprop}{Cómo se resuelven}
\begin{itemize}
\item Un \textbf{sistema de inecuaciones} con una incógnita es un conjunto de inecuaciones que se tienen que cumplir \textbf{a la vez}.
\item Se resuelve cada inecuación por separado y la solución del sistema es la \textbf{intersección} (la parte común) de todas las soluciones.
\item Conviene representar cada solución en una recta real, una sobre la otra y con la misma escala. Si no tienen ningún punto en común, el sistema \textbf{no tiene solución}.
\end{itemize}
\end{fmprop}

\Needspace{5\baselineskip}
\begin{fmex}{Ejemplo resuelto: sistemas de inecuaciones}
\begin{enumerate}[label=\textbf{\arabic*.},leftmargin=*]
\item Resuelve el sistema \fmm{\begin{cases}3x-2\le x+4\\ 2x+1\gt-3\end{cases}}

\fmm{2x\le6\Rightarrow x\le3:\ (-\infty,3].\qquad 2x\gt-4\Rightarrow x\gt-2:\ (-2,+\infty)}

\begin{center}\begin{tikzpicture}[line cap=round]\draw[fmmuted!60,dashed,line width=0.5pt] (4.678,0.1)--(4.678,-1.700);\draw[fmmuted!60,dashed,line width=0.5pt] (10.622,0.1)--(10.622,-1.700);\node[anchor=west,font=\small] at (0,-0.400) {$3x-2\leq x+4$};\draw[fmmuted!50,line width=0.6pt] (2.300,-0.400)--(13.000,-0.400);\draw[fmA,line width=2.6pt] (2.300,-0.400)--(10.622,-0.400);\fill[fmA] (2.050,-0.400)--++(0.3,0.19)--++(0,-0.38)--cycle;\draw[fmA,line width=1.2pt,fill=fmA] (10.622,-0.400) circle[radius=0.13];\node[anchor=west,font=\small] at (0,-1.150) {$2x+1>-3$};\draw[fmmuted!50,line width=0.6pt] (2.300,-1.150)--(13.000,-1.150);\draw[fmB,line width=2.6pt] (4.678,-1.150)--(13.000,-1.150);\fill[fmB] (13.250,-1.150)--++(-0.3,0.19)--++(0,-0.38)--cycle;\draw[fmB,line width=1.2pt,fill=white] (4.678,-1.150) circle[radius=0.13];\draw[fmmuted,line width=0.8pt,-latex] (2.100,-1.700)--(13.400,-1.700);\draw[fmmuted,line width=0.6pt] (2.300,-1.700)--++(0,0.12)--++(0,-0.24) node[below,font=\scriptsize,text=fmmuted] {$-4$};\draw[fmmuted,line width=0.6pt] (3.489,-1.700)--++(0,0.12)--++(0,-0.24) node[below,font=\scriptsize,text=fmmuted] {$-3$};\draw[fmmuted,line width=0.6pt] (4.678,-1.700)--++(0,0.12)--++(0,-0.24) node[below,font=\scriptsize,text=fmmuted] {$-2$};\draw[fmmuted,line width=0.6pt] (5.867,-1.700)--++(0,0.12)--++(0,-0.24) node[below,font=\scriptsize,text=fmmuted] {$-1$};\draw[fmmuted,line width=0.6pt] (7.056,-1.700)--++(0,0.12)--++(0,-0.24) node[below,font=\scriptsize,text=fmmuted] {$0$};\draw[fmmuted,line width=0.6pt] (8.244,-1.700)--++(0,0.12)--++(0,-0.24) node[below,font=\scriptsize,text=fmmuted] {$1$};\draw[fmmuted,line width=0.6pt] (9.433,-1.700)--++(0,0.12)--++(0,-0.24) node[below,font=\scriptsize,text=fmmuted] {$2$};\draw[fmmuted,line width=0.6pt] (10.622,-1.700)--++(0,0.12)--++(0,-0.24) node[below,font=\scriptsize,text=fmmuted] {$3$};\draw[fmmuted,line width=0.6pt] (11.811,-1.700)--++(0,0.12)--++(0,-0.24) node[below,font=\scriptsize,text=fmmuted] {$4$};\draw[fmmuted,line width=0.6pt] (13.000,-1.700)--++(0,0.12)--++(0,-0.24) node[below,font=\scriptsize,text=fmmuted] {$5$};\end{tikzpicture}\end{center}

La solución es la parte donde coinciden las dos rectas: \fmm{x\in(-2,\,3]}

\begin{center}\begin{tikzpicture}[line cap=round]\fill[fmR,opacity=0.12] (4.678,0.1) rectangle (10.622,-2.450);\draw[fmmuted!60,dashed,line width=0.5pt] (4.678,0.1)--(4.678,-2.450);\draw[fmmuted!60,dashed,line width=0.5pt] (10.622,0.1)--(10.622,-2.450);\node[anchor=west,font=\small] at (0,-0.400) {$3x-2\leq x+4$};\draw[fmmuted!50,line width=0.6pt] (2.300,-0.400)--(13.000,-0.400);\draw[fmA,line width=2.6pt] (2.300,-0.400)--(10.622,-0.400);\fill[fmA] (2.050,-0.400)--++(0.3,0.19)--++(0,-0.38)--cycle;\draw[fmA,line width=1.2pt,fill=fmA] (10.622,-0.400) circle[radius=0.13];\node[anchor=west,font=\small] at (0,-1.150) {$2x+1>-3$};\draw[fmmuted!50,line width=0.6pt] (2.300,-1.150)--(13.000,-1.150);\draw[fmB,line width=2.6pt] (4.678,-1.150)--(13.000,-1.150);\fill[fmB] (13.250,-1.150)--++(-0.3,0.19)--++(0,-0.38)--cycle;\draw[fmB,line width=1.2pt,fill=white] (4.678,-1.150) circle[radius=0.13];\node[anchor=west,font=\small] at (0,-1.900) {$\text{Solución}$};\draw[fmmuted!50,line width=0.6pt] (2.300,-1.900)--(13.000,-1.900);\draw[fmR,line width=2.6pt] (4.678,-1.900)--(10.622,-1.900);\draw[fmR,line width=1.2pt,fill=white] (4.678,-1.900) circle[radius=0.13];\draw[fmR,line width=1.2pt,fill=fmR] (10.622,-1.900) circle[radius=0.13];\draw[fmmuted,line width=0.8pt,-latex] (2.100,-2.450)--(13.400,-2.450);\draw[fmmuted,line width=0.6pt] (2.300,-2.450)--++(0,0.12)--++(0,-0.24) node[below,font=\scriptsize,text=fmmuted] {$-4$};\draw[fmmuted,line width=0.6pt] (3.489,-2.450)--++(0,0.12)--++(0,-0.24) node[below,font=\scriptsize,text=fmmuted] {$-3$};\draw[fmmuted,line width=0.6pt] (4.678,-2.450)--++(0,0.12)--++(0,-0.24) node[below,font=\scriptsize,text=fmmuted] {$-2$};\draw[fmmuted,line width=0.6pt] (5.867,-2.450)--++(0,0.12)--++(0,-0.24) node[below,font=\scriptsize,text=fmmuted] {$-1$};\draw[fmmuted,line width=0.6pt] (7.056,-2.450)--++(0,0.12)--++(0,-0.24) node[below,font=\scriptsize,text=fmmuted] {$0$};\draw[fmmuted,line width=0.6pt] (8.244,-2.450)--++(0,0.12)--++(0,-0.24) node[below,font=\scriptsize,text=fmmuted] {$1$};\draw[fmmuted,line width=0.6pt] (9.433,-2.450)--++(0,0.12)--++(0,-0.24) node[below,font=\scriptsize,text=fmmuted] {$2$};\draw[fmmuted,line width=0.6pt] (10.622,-2.450)--++(0,0.12)--++(0,-0.24) node[below,font=\scriptsize,text=fmmuted] {$3$};\draw[fmmuted,line width=0.6pt] (11.811,-2.450)--++(0,0.12)--++(0,-0.24) node[below,font=\scriptsize,text=fmmuted] {$4$};\draw[fmmuted,line width=0.6pt] (13.000,-2.450)--++(0,0.12)--++(0,-0.24) node[below,font=\scriptsize,text=fmmuted] {$5$};\end{tikzpicture}\end{center}
\item Resuelve el sistema \fmm{\begin{cases}x^2-x-6\le0\\ x-1\gt0\end{cases}}

Raíces de \fmm{x^2-x-6}: \fmm{-2} y \fmm{3}, y es negativo entre ellas: \fmm{[-2,3]}. \fmm{\quad x\gt1:\ (1,+\infty)}

\begin{center}\begin{tikzpicture}[line cap=round]\draw[fmmuted!60,dashed,line width=0.5pt] (4.678,0.1)--(4.678,-1.700);\draw[fmmuted!60,dashed,line width=0.5pt] (8.244,0.1)--(8.244,-1.700);\draw[fmmuted!60,dashed,line width=0.5pt] (10.622,0.1)--(10.622,-1.700);\node[anchor=west,font=\small] at (0,-0.400) {$x^2-x-6\leq 0$};\draw[fmmuted!50,line width=0.6pt] (2.300,-0.400)--(13.000,-0.400);\draw[fmA,line width=2.6pt] (4.678,-0.400)--(10.622,-0.400);\draw[fmA,line width=1.2pt,fill=fmA] (4.678,-0.400) circle[radius=0.13];\draw[fmA,line width=1.2pt,fill=fmA] (10.622,-0.400) circle[radius=0.13];\node[anchor=west,font=\small] at (0,-1.150) {$x-1>0$};\draw[fmmuted!50,line width=0.6pt] (2.300,-1.150)--(13.000,-1.150);\draw[fmB,line width=2.6pt] (8.244,-1.150)--(13.000,-1.150);\fill[fmB] (13.250,-1.150)--++(-0.3,0.19)--++(0,-0.38)--cycle;\draw[fmB,line width=1.2pt,fill=white] (8.244,-1.150) circle[radius=0.13];\draw[fmmuted,line width=0.8pt,-latex] (2.100,-1.700)--(13.400,-1.700);\draw[fmmuted,line width=0.6pt] (2.300,-1.700)--++(0,0.12)--++(0,-0.24) node[below,font=\scriptsize,text=fmmuted] {$-4$};\draw[fmmuted,line width=0.6pt] (3.489,-1.700)--++(0,0.12)--++(0,-0.24) node[below,font=\scriptsize,text=fmmuted] {$-3$};\draw[fmmuted,line width=0.6pt] (4.678,-1.700)--++(0,0.12)--++(0,-0.24) node[below,font=\scriptsize,text=fmmuted] {$-2$};\draw[fmmuted,line width=0.6pt] (5.867,-1.700)--++(0,0.12)--++(0,-0.24) node[below,font=\scriptsize,text=fmmuted] {$-1$};\draw[fmmuted,line width=0.6pt] (7.056,-1.700)--++(0,0.12)--++(0,-0.24) node[below,font=\scriptsize,text=fmmuted] {$0$};\draw[fmmuted,line width=0.6pt] (8.244,-1.700)--++(0,0.12)--++(0,-0.24) node[below,font=\scriptsize,text=fmmuted] {$1$};\draw[fmmuted,line width=0.6pt] (9.433,-1.700)--++(0,0.12)--++(0,-0.24) node[below,font=\scriptsize,text=fmmuted] {$2$};\draw[fmmuted,line width=0.6pt] (10.622,-1.700)--++(0,0.12)--++(0,-0.24) node[below,font=\scriptsize,text=fmmuted] {$3$};\draw[fmmuted,line width=0.6pt] (11.811,-1.700)--++(0,0.12)--++(0,-0.24) node[below,font=\scriptsize,text=fmmuted] {$4$};\draw[fmmuted,line width=0.6pt] (13.000,-1.700)--++(0,0.12)--++(0,-0.24) node[below,font=\scriptsize,text=fmmuted] {$5$};\end{tikzpicture}\end{center}

\fmm{x\in(1,\,3]}. El 1 no se incluye porque no cumple la segunda inecuación.

\begin{center}\begin{tikzpicture}[line cap=round]\fill[fmR,opacity=0.12] (8.244,0.1) rectangle (10.622,-2.450);\draw[fmmuted!60,dashed,line width=0.5pt] (4.678,0.1)--(4.678,-2.450);\draw[fmmuted!60,dashed,line width=0.5pt] (8.244,0.1)--(8.244,-2.450);\draw[fmmuted!60,dashed,line width=0.5pt] (10.622,0.1)--(10.622,-2.450);\node[anchor=west,font=\small] at (0,-0.400) {$x^2-x-6\leq 0$};\draw[fmmuted!50,line width=0.6pt] (2.300,-0.400)--(13.000,-0.400);\draw[fmA,line width=2.6pt] (4.678,-0.400)--(10.622,-0.400);\draw[fmA,line width=1.2pt,fill=fmA] (4.678,-0.400) circle[radius=0.13];\draw[fmA,line width=1.2pt,fill=fmA] (10.622,-0.400) circle[radius=0.13];\node[anchor=west,font=\small] at (0,-1.150) {$x-1>0$};\draw[fmmuted!50,line width=0.6pt] (2.300,-1.150)--(13.000,-1.150);\draw[fmB,line width=2.6pt] (8.244,-1.150)--(13.000,-1.150);\fill[fmB] (13.250,-1.150)--++(-0.3,0.19)--++(0,-0.38)--cycle;\draw[fmB,line width=1.2pt,fill=white] (8.244,-1.150) circle[radius=0.13];\node[anchor=west,font=\small] at (0,-1.900) {$\text{Solución}$};\draw[fmmuted!50,line width=0.6pt] (2.300,-1.900)--(13.000,-1.900);\draw[fmR,line width=2.6pt] (8.244,-1.900)--(10.622,-1.900);\draw[fmR,line width=1.2pt,fill=white] (8.244,-1.900) circle[radius=0.13];\draw[fmR,line width=1.2pt,fill=fmR] (10.622,-1.900) circle[radius=0.13];\draw[fmmuted,line width=0.8pt,-latex] (2.100,-2.450)--(13.400,-2.450);\draw[fmmuted,line width=0.6pt] (2.300,-2.450)--++(0,0.12)--++(0,-0.24) node[below,font=\scriptsize,text=fmmuted] {$-4$};\draw[fmmuted,line width=0.6pt] (3.489,-2.450)--++(0,0.12)--++(0,-0.24) node[below,font=\scriptsize,text=fmmuted] {$-3$};\draw[fmmuted,line width=0.6pt] (4.678,-2.450)--++(0,0.12)--++(0,-0.24) node[below,font=\scriptsize,text=fmmuted] {$-2$};\draw[fmmuted,line width=0.6pt] (5.867,-2.450)--++(0,0.12)--++(0,-0.24) node[below,font=\scriptsize,text=fmmuted] {$-1$};\draw[fmmuted,line width=0.6pt] (7.056,-2.450)--++(0,0.12)--++(0,-0.24) node[below,font=\scriptsize,text=fmmuted] {$0$};\draw[fmmuted,line width=0.6pt] (8.244,-2.450)--++(0,0.12)--++(0,-0.24) node[below,font=\scriptsize,text=fmmuted] {$1$};\draw[fmmuted,line width=0.6pt] (9.433,-2.450)--++(0,0.12)--++(0,-0.24) node[below,font=\scriptsize,text=fmmuted] {$2$};\draw[fmmuted,line width=0.6pt] (10.622,-2.450)--++(0,0.12)--++(0,-0.24) node[below,font=\scriptsize,text=fmmuted] {$3$};\draw[fmmuted,line width=0.6pt] (11.811,-2.450)--++(0,0.12)--++(0,-0.24) node[below,font=\scriptsize,text=fmmuted] {$4$};\draw[fmmuted,line width=0.6pt] (13.000,-2.450)--++(0,0.12)--++(0,-0.24) node[below,font=\scriptsize,text=fmmuted] {$5$};\end{tikzpicture}\end{center}
\item Resuelve el sistema \fmm{\begin{cases}x-2\gt0\\ x+1\lt0\end{cases}}

\fmm{x\gt2} y \fmm{x\lt-1}.

\begin{center}\begin{tikzpicture}[line cap=round]\draw[fmmuted!60,dashed,line width=0.5pt] (5.357,0.1)--(5.357,-1.700);\draw[fmmuted!60,dashed,line width=0.5pt] (9.943,0.1)--(9.943,-1.700);\node[anchor=west,font=\small] at (0,-0.400) {$x-2>0$};\draw[fmmuted!50,line width=0.6pt] (2.300,-0.400)--(13.000,-0.400);\draw[fmA,line width=2.6pt] (9.943,-0.400)--(13.000,-0.400);\fill[fmA] (13.250,-0.400)--++(-0.3,0.19)--++(0,-0.38)--cycle;\draw[fmA,line width=1.2pt,fill=white] (9.943,-0.400) circle[radius=0.13];\node[anchor=west,font=\small] at (0,-1.150) {$x+1<0$};\draw[fmmuted!50,line width=0.6pt] (2.300,-1.150)--(13.000,-1.150);\draw[fmB,line width=2.6pt] (2.300,-1.150)--(5.357,-1.150);\fill[fmB] (2.050,-1.150)--++(0.3,0.19)--++(0,-0.38)--cycle;\draw[fmB,line width=1.2pt,fill=white] (5.357,-1.150) circle[radius=0.13];\draw[fmmuted,line width=0.8pt,-latex] (2.100,-1.700)--(13.400,-1.700);\draw[fmmuted,line width=0.6pt] (2.300,-1.700)--++(0,0.12)--++(0,-0.24) node[below,font=\scriptsize,text=fmmuted] {$-3$};\draw[fmmuted,line width=0.6pt] (3.829,-1.700)--++(0,0.12)--++(0,-0.24) node[below,font=\scriptsize,text=fmmuted] {$-2$};\draw[fmmuted,line width=0.6pt] (5.357,-1.700)--++(0,0.12)--++(0,-0.24) node[below,font=\scriptsize,text=fmmuted] {$-1$};\draw[fmmuted,line width=0.6pt] (6.886,-1.700)--++(0,0.12)--++(0,-0.24) node[below,font=\scriptsize,text=fmmuted] {$0$};\draw[fmmuted,line width=0.6pt] (8.414,-1.700)--++(0,0.12)--++(0,-0.24) node[below,font=\scriptsize,text=fmmuted] {$1$};\draw[fmmuted,line width=0.6pt] (9.943,-1.700)--++(0,0.12)--++(0,-0.24) node[below,font=\scriptsize,text=fmmuted] {$2$};\draw[fmmuted,line width=0.6pt] (11.471,-1.700)--++(0,0.12)--++(0,-0.24) node[below,font=\scriptsize,text=fmmuted] {$3$};\draw[fmmuted,line width=0.6pt] (13.000,-1.700)--++(0,0.12)--++(0,-0.24) node[below,font=\scriptsize,text=fmmuted] {$4$};\end{tikzpicture}\end{center}

No hay ningún número que sea mayor que 2 y menor que \fmm{-1}: el sistema no tiene solución.

\begin{center}\begin{tikzpicture}[line cap=round]\draw[fmmuted!60,dashed,line width=0.5pt] (5.357,0.1)--(5.357,-2.450);\draw[fmmuted!60,dashed,line width=0.5pt] (9.943,0.1)--(9.943,-2.450);\node[anchor=west,font=\small] at (0,-0.400) {$x-2>0$};\draw[fmmuted!50,line width=0.6pt] (2.300,-0.400)--(13.000,-0.400);\draw[fmA,line width=2.6pt] (9.943,-0.400)--(13.000,-0.400);\fill[fmA] (13.250,-0.400)--++(-0.3,0.19)--++(0,-0.38)--cycle;\draw[fmA,line width=1.2pt,fill=white] (9.943,-0.400) circle[radius=0.13];\node[anchor=west,font=\small] at (0,-1.150) {$x+1<0$};\draw[fmmuted!50,line width=0.6pt] (2.300,-1.150)--(13.000,-1.150);\draw[fmB,line width=2.6pt] (2.300,-1.150)--(5.357,-1.150);\fill[fmB] (2.050,-1.150)--++(0.3,0.19)--++(0,-0.38)--cycle;\draw[fmB,line width=1.2pt,fill=white] (5.357,-1.150) circle[radius=0.13];\node[anchor=west,font=\small] at (0,-1.900) {$\text{Solución}$};\draw[fmmuted!50,line width=0.6pt] (2.300,-1.900)--(13.000,-1.900);\node[font=\small] at (7.650,-1.650) {$\emptyset$};\draw[fmmuted,line width=0.8pt,-latex] (2.100,-2.450)--(13.400,-2.450);\draw[fmmuted,line width=0.6pt] (2.300,-2.450)--++(0,0.12)--++(0,-0.24) node[below,font=\scriptsize,text=fmmuted] {$-3$};\draw[fmmuted,line width=0.6pt] (3.829,-2.450)--++(0,0.12)--++(0,-0.24) node[below,font=\scriptsize,text=fmmuted] {$-2$};\draw[fmmuted,line width=0.6pt] (5.357,-2.450)--++(0,0.12)--++(0,-0.24) node[below,font=\scriptsize,text=fmmuted] {$-1$};\draw[fmmuted,line width=0.6pt] (6.886,-2.450)--++(0,0.12)--++(0,-0.24) node[below,font=\scriptsize,text=fmmuted] {$0$};\draw[fmmuted,line width=0.6pt] (8.414,-2.450)--++(0,0.12)--++(0,-0.24) node[below,font=\scriptsize,text=fmmuted] {$1$};\draw[fmmuted,line width=0.6pt] (9.943,-2.450)--++(0,0.12)--++(0,-0.24) node[below,font=\scriptsize,text=fmmuted] {$2$};\draw[fmmuted,line width=0.6pt] (11.471,-2.450)--++(0,0.12)--++(0,-0.24) node[below,font=\scriptsize,text=fmmuted] {$3$};\draw[fmmuted,line width=0.6pt] (13.000,-2.450)--++(0,0.12)--++(0,-0.24) node[below,font=\scriptsize,text=fmmuted] {$4$};\end{tikzpicture}\end{center}
\item Buscamos los números enteros tales que su doble menos 3 es menor que 11 y su triple más 1 es mayor que 7.

\fmm{\begin{cases}2x-3\lt11\\ 3x+1\gt7\end{cases}\Rightarrow\begin{cases}x\lt7\\ x\gt2\end{cases}\Rightarrow x\in(2,\,7)}

Los enteros del intervalo son 3, 4, 5 y 6.
\end{enumerate}
\end{fmex}

\Needspace{5.7cm}
\section{Inecuaciones con dos incógnitas}\hypertarget{sec7}{}
\begin{fmdef}{Inecuaciones lineales con dos incógnitas}
Una inecuación lineal con dos incógnitas tiene la forma \fmm{ax+by\lt c} (o con \fmm{\gt,\le,\ge}). Una \textbf{solución} es un par de números \fmm{(x,y)} que la cumple, es decir, un punto del plano. Hay infinitas: todas juntas forman un \textbf{semiplano}, una de las dos regiones en que la recta \fmm{ax+by=c} divide el plano.

\begin{fmnote}Aparecen en muchos problemas de la economía. Por ejemplo, si con 30 € como máximo compramos \fmm{x} entradas de 5 € e \fmm{y} entradas de 3 €, debe ser \fmm{5x+3y\le30}.\end{fmnote}
\end{fmdef}

\Needspace{5\baselineskip}
\begin{fmprop}{Cómo se resuelven, paso a paso}
\begin{enumerate}
\item \textbf{Recta frontera.} Cambiamos la desigualdad por un igual: \fmm{ax+by=c}.
\item \textbf{Dibujamos la recta} a partir de dos puntos (por ejemplo, los cortes con los ejes, haciendo \fmm{x=0} e \fmm{y=0}). Es \textbf{continua} si la desigualdad es \fmm{\le} o \fmm{\ge}, porque los puntos de la recta son solución, y \textbf{discontinua} si es \fmm{\lt} o \fmm{\gt}.
\item \textbf{Punto de prueba.} Elegimos un punto que no esté en la recta (a menudo \fmm{O(0,0)}; si la recta pasa por el origen, otro, como \fmm{(1,0)}) y lo sustituimos en la inecuación.
\item \textbf{Elegimos el semiplano.} Si el punto cumple la inecuación, la solución es el semiplano que lo contiene; si no, es el otro. Lo sombreamos.
\end{enumerate}
\end{fmprop}

\Needspace{5\baselineskip}
\begin{fmex}{Ejemplo resuelto: inecuaciones con dos incógnitas}
\begin{enumerate}[label=\textbf{\arabic*.},leftmargin=*]
\item Resuelve \fmm{x+2y\le6}.

Recta frontera: \fmm{x+2y=6}. Dos puntos: \fmm{(6,0)} y \fmm{(0,3)}. Como hay \fmm{\le}, la recta es continua.

Punto de prueba \fmm{O(0,0)}: \fmm{0+2\cdot0=0\le6\ \checkmark}

La solución es el semiplano que contiene el origen, con la recta incluida.

\begin{center}\begin{tikzpicture}\begin{axis}[fmaxis,axis equal image,width=4.6cm,height=4.6cm,xmin=-3,xmax=7,ymin=-3,ymax=7,xtick={-2,2,4,6},ytick={-2,2,4,6},minor xtick={-3,-2,-1,0,1,2,3,4,5,6,7},minor ytick={-3,-2,-1,0,1,2,3,4,5,6,7},grid=both,minor grid style={fmmuted!15},major grid style={fmmuted!15},tick label style={font=\tiny}]\fill[fmR,opacity=0.22] (axis cs:-3.0000,-3.0000)--(axis cs:7.0000,-3.0000)--(axis cs:7.0000,-0.5000)--(axis cs:-3.0000,4.5000)--cycle;\draw[fmA,line width=1.5pt] (axis cs:-4.0000,5.0000)--(axis cs:8.0000,-1.0000);\node[circle,fill=fmK,inner sep=0pt,minimum size=6pt] at (axis cs:0,0) {};\end{axis}\end{tikzpicture}\end{center}
\item Resuelve \fmm{y\gt x+1}.

Recta frontera: \fmm{y=x+1}. Dos puntos: \fmm{(-1,0)} y \fmm{(0,1)}. Como hay \fmm{\gt}, la recta es discontinua.

Punto de prueba \fmm{O(0,0)}: \fmm{0\gt0+1} es falso.

La solución es el semiplano que no contiene el origen, sin la recta.

\begin{center}\begin{tikzpicture}\begin{axis}[fmaxis,axis equal image,width=4.6cm,height=4.6cm,xmin=-5,xmax=5,ymin=-5,ymax=5,xtick={-4,-2,2,4},ytick={-4,-2,2,4},minor xtick={-5,-4,-3,-2,-1,0,1,2,3,4,5},minor ytick={-5,-4,-3,-2,-1,0,1,2,3,4,5},grid=both,minor grid style={fmmuted!15},major grid style={fmmuted!15},tick label style={font=\tiny}]\fill[fmR,opacity=0.22] (axis cs:4.0000,5.0000)--(axis cs:-5.0000,5.0000)--(axis cs:-5.0000,-4.0000)--cycle;\draw[fmA,line width=1.5pt,dash pattern=on 5pt off 3.5pt] (axis cs:-6.0000,-5.0000)--(axis cs:6.0000,7.0000);\node[circle,fill=fmR,inner sep=0pt,minimum size=6pt] at (axis cs:0,0) {};\end{axis}\end{tikzpicture}\end{center}
\item Resuelve \fmm{y\ge2x}.

Recta frontera: \fmm{y=2x}. Dos puntos: \fmm{(0,0)} y \fmm{(1,2)}. Recta continua.

La recta pasa por el origen, así que probamos otro punto: \fmm{(1,0)}: \fmm{0\ge2\cdot1} es falso.

La solución es el semiplano que no contiene el punto \fmm{(1,0)}, con la recta incluida.

\begin{center}\begin{tikzpicture}\begin{axis}[fmaxis,axis equal image,width=4.6cm,height=4.6cm,xmin=-5,xmax=5,ymin=-5,ymax=5,xtick={-4,-2,2,4},ytick={-4,-2,2,4},minor xtick={-5,-4,-3,-2,-1,0,1,2,3,4,5},minor ytick={-5,-4,-3,-2,-1,0,1,2,3,4,5},grid=both,minor grid style={fmmuted!15},major grid style={fmmuted!15},tick label style={font=\tiny}]\fill[fmR,opacity=0.22] (axis cs:-5.0000,-5.0000)--(axis cs:-2.5000,-5.0000)--(axis cs:2.5000,5.0000)--(axis cs:-5.0000,5.0000)--cycle;\draw[fmA,line width=1.5pt] (axis cs:-6.0000,-12.0000)--(axis cs:6.0000,12.0000);\node[circle,fill=fmR,inner sep=0pt,minimum size=6pt] at (axis cs:1,0) {};\end{axis}\end{tikzpicture}\end{center}
\end{enumerate}
\end{fmex}

\Needspace{5.7cm}
\section{Sistemas de inecuaciones con dos incógnitas}\hypertarget{sec8}{}
\begin{fmprop}{Cómo se resuelven}
\begin{enumerate}
\item Se resuelve cada inecuación por separado, como en el apartado anterior: cada una da un semiplano.
\item La solución del sistema es la \textbf{intersección} de todos los semiplanos, la \textbf{región factible}. Si no tienen ningún punto en común, el sistema no tiene solución.
\item Los \textbf{vértices} de la región son los puntos de corte de las rectas frontera: se calculan resolviendo el sistema formado por las dos ecuaciones.
\item La región puede ser \textbf{acotada} (un polígono) o \textbf{no acotada} (se extiende indefinidamente). Las restricciones \fmm{x\ge0,\ y\ge0} la limitan al primer cuadrante.
\end{enumerate}
\end{fmprop}

\Needspace{5\baselineskip}
\begin{fmex}{Ejemplo resuelto: sistemas con dos incógnitas}
\begin{enumerate}[label=\textbf{\arabic*.},leftmargin=*]
\item Resuelve \fmm{\begin{cases}x+y\le5\\ x-y\ge-1\\ x\ge0,\ y\ge0\end{cases}}

\fmm{x+y\le5}: recta \fmm{x+y=5}; con \fmm{O(0,0)}: \fmm{0\le5\ \checkmark}: semiplano de \fmm{O}.

\fmm{x-y\ge-1}: recta \fmm{x-y=-1}, por \fmm{(0,1)} y \fmm{(-1,0)}; con \fmm{O}: \fmm{0\ge-1\ \checkmark}: semiplano de \fmm{O}. \fmm{\ x\ge0,\ y\ge0}: primer cuadrante.

Vértice donde se cortan las dos rectas oblicuas: \fmm{\begin{cases}x+y=5\\ x-y=-1\end{cases}\Rightarrow 2x=4,\ x=2,\ y=3}

La solución es el cuadrilátero de vértices \fmm{(0,0),\ (5,0),\ (2,3)} y \fmm{(0,1)}. La región es acotada.

\begin{center}\begin{tikzpicture}\begin{axis}[fmaxis,axis equal image,width=4.6cm,height=4.6cm,xmin=-2,xmax=6,ymin=-2,ymax=6,xtick={-2,2,4,6},ytick={-2,2,4,6},minor xtick={-2,-1,0,1,2,3,4,5,6},minor ytick={-2,-1,0,1,2,3,4,5,6},grid=both,minor grid style={fmmuted!15},major grid style={fmmuted!15},tick label style={font=\tiny}]\fill[fmR,opacity=0.22] (axis cs:5.0000,0.0000)--(axis cs:2.0000,3.0000)--(axis cs:0.0000,1.0000)--(axis cs:0.0000,0.0000)--cycle;\draw[fmA,line width=1.5pt] (axis cs:-3.0000,8.0000)--(axis cs:7.0000,-2.0000);\draw[fmB,line width=1.5pt] (axis cs:-3.0000,-2.0000)--(axis cs:7.0000,8.0000);\draw[fmV,line width=1.5pt] (axis cs:0.0000,-3.0000)--(axis cs:0.0000,7.0000);\draw[fmK,line width=1.5pt] (axis cs:-3.0000,0.0000)--(axis cs:7.0000,0.0000);\end{axis}\end{tikzpicture}\end{center}
\item Resuelve \fmm{\begin{cases}x+y\ge2\\ x-y\le2\end{cases}}

Con \fmm{O}: \fmm{0\ge2} es falso (semiplano sin \fmm{O}); \fmm{0\le2\ \checkmark} (semiplano de \fmm{O}). Las rectas se cortan en \fmm{(2,0)}.

Región \textbf{no acotada} con un solo vértice, \fmm{(2,0)}.

\begin{center}\begin{tikzpicture}\begin{axis}[fmaxis,axis equal image,width=4.6cm,height=4.6cm,xmin=-3,xmax=6,ymin=-3,ymax=6,xtick={-2,2,4,6},ytick={-2,2,4,6},minor xtick={-3,-2,-1,0,1,2,3,4,5,6},minor ytick={-3,-2,-1,0,1,2,3,4,5,6},grid=both,minor grid style={fmmuted!15},major grid style={fmmuted!15},tick label style={font=\tiny}]\fill[fmR,opacity=0.22] (axis cs:6.0000,4.0000)--(axis cs:6.0000,6.0000)--(axis cs:-3.0000,6.0000)--(axis cs:-3.0000,5.0000)--(axis cs:2.0000,0.0000)--cycle;\draw[fmA,line width=1.5pt] (axis cs:-4.0000,6.0000)--(axis cs:7.0000,-5.0000);\draw[fmB,line width=1.5pt] (axis cs:-4.0000,-6.0000)--(axis cs:7.0000,5.0000);\end{axis}\end{tikzpicture}\end{center}
\item Resuelve \fmm{\begin{cases}x+y\le1\\ x+y\ge3\end{cases}}

Las rectas \fmm{x+y=1} y \fmm{x+y=3} son paralelas y los semiplanos no tienen ningún punto en común (ningún par de números puede sumar a la vez como máximo 1 y como mínimo 3): el sistema no tiene solución.

\begin{center}\begin{tikzpicture}\begin{axis}[fmaxis,axis equal image,width=4.6cm,height=4.6cm,xmin=-2,xmax=5,ymin=-2,ymax=5,xtick={-2,2,4},ytick={-2,2,4},minor xtick={-2,-1,0,1,2,3,4,5},minor ytick={-2,-1,0,1,2,3,4,5},grid=both,minor grid style={fmmuted!15},major grid style={fmmuted!15},tick label style={font=\tiny}]\draw[fmA,line width=1.5pt] (axis cs:-3.0000,4.0000)--(axis cs:6.0000,-5.0000);\draw[fmB,line width=1.5pt] (axis cs:-3.0000,6.0000)--(axis cs:6.0000,-3.0000);\end{axis}\end{tikzpicture}\end{center}
\end{enumerate}
\end{fmex}

\Needspace{5.7cm}
\section{Programación lineal}\hypertarget{sec9}{}
\begin{fmdef}{Qué es un problema de programación lineal}
Muchas decisiones de empresas y familias consisten en hacer \textbf{máximo} un beneficio o \textbf{mínimo} un coste con unos recursos limitados. Un problema de \textbf{programación lineal} con dos variables tiene:

\begin{itemize}
\item Una \textbf{función objetivo} \fmm{F(x,y)=ax+by}, que queremos maximizar o minimizar.
\item Unas \textbf{restricciones}: un sistema de inecuaciones lineales (normalmente con \fmm{x\ge0,\ y\ge0}).
\item La \textbf{región factible}: el conjunto de puntos que cumplen todas las restricciones (las soluciones posibles). La \textbf{solución óptima} es el punto de la región donde \fmm{F} toma el valor máximo o mínimo.
\end{itemize}
\end{fmdef}

\Needspace{5\baselineskip}
\begin{fmprop}{Teorema fundamental}
Si la región factible es acotada, la función objetivo alcanza el máximo y el mínimo en un \textbf{vértice} de la región. Si los alcanza en dos vértices, también los alcanza en todos los puntos del lado que los une (\textbf{infinitas soluciones}). Por tanto, basta con calcular \fmm{F} en los vértices.

\begin{fmnote}Si la región no es acotada, puede que no haya máximo (o mínimo): hay que mirar hacia dónde crece \fmm{F}. Si la región es vacía, el problema no tiene solución.\end{fmnote}
\end{fmprop}

\Needspace{5\baselineskip}
\begin{fmprop}{Pasos para resolver un problema}
\begin{enumerate}
\item \textbf{Nombramos las incógnitas} \fmm{x} e \fmm{y} y, si ayuda, organizamos los datos en una \textbf{tabla}.
\item Escribimos las \textbf{restricciones} («como máximo» es \fmm{\le}; «como mínimo» es \fmm{\ge}) y la \textbf{función objetivo}.
\item Dibujamos la \textbf{región factible} y calculamos sus \textbf{vértices}.
\item \textbf{Evaluamos} \fmm{F} en cada vértice y elegimos el valor más grande (máximo) o el más pequeño (mínimo).
\item \textbf{Respondemos} la pregunta con una frase.
\end{enumerate}
\end{fmprop}

\Needspace{5\baselineskip}
\begin{fmex}{Ejemplo resuelto: un problema de máximo}
\begin{enumerate}[label=\textbf{\arabic*.},leftmargin=*]
\item Una fábrica de juguetes hace coches y trenes. Cada coche necesita 1 hora de montaje y 1 hora de acabado, y cada tren, 2 horas de montaje y 1 de acabado. Cada día hay como máximo 16 horas de montaje y 10 de acabado. Cada coche da un beneficio de 30 € y cada tren, de 40 €. ¿Cuántos tiene que hacer de cada tipo para obtener el beneficio máximo?

\fmm{x=} coches, \fmm{y=} trenes.

\begin{center}\renewcommand{\arraystretch}{1.5}\setlength{\parskip}{0pt}\small\begin{tabularx}{\linewidth}{>{\raggedright\arraybackslash}p{3.2cm}>{\raggedright\arraybackslash}X>{\raggedright\arraybackslash}X>{\raggedright\arraybackslash}X}\toprule
\textbf{} & \textbf{Montaje (h)} & \textbf{Acabado (h)} & \textbf{Beneficio (€)}\\\midrule
Coches $(x)$ & 1 & 1 & 30\\[3pt]\end{tabularx}\par\nointerlineskip\vskip0pt
\small\begin{tabularx}{\linewidth}{>{\raggedright\arraybackslash}p{3.2cm}>{\raggedright\arraybackslash}X>{\raggedright\arraybackslash}X>{\raggedright\arraybackslash}X}Trenes $(y)$ & 2 & 1 & 40\\[3pt]\end{tabularx}\par\nointerlineskip\vskip0pt
\small\begin{tabularx}{\linewidth}{>{\raggedright\arraybackslash}p{3.2cm}>{\raggedright\arraybackslash}X>{\raggedright\arraybackslash}X>{\raggedright\arraybackslash}X}Disponible & 16 & 10 & \\[3pt]\bottomrule\end{tabularx}\end{center}

Restricciones: \fmm{\begin{cases}x+2y\le16\\ x+y\le10\\ x\ge0,\ y\ge0\end{cases}\qquad} Función objetivo: \fmm{F(x,y)=30x+40y} (máximo)

Región factible y vértices. El vértice de arriba es el corte de \fmm{x+2y=16} y \fmm{x+y=10}: restando, \fmm{y=6} y \fmm{x=4}. \fmm{\ A(0,0),\ B(10,0),\ C(4,6),\ D(0,8)}

\begin{center}\begin{tikzpicture}\begin{axis}[fmaxis,axis equal image,width=4.6cm,height=4.6cm,xmin=-1,xmax=17,ymin=-1,ymax=17,xtick={2,4,6,8,10,12,14,16},ytick={2,4,6,8,10,12,14,16},minor xtick={-1,0,1,2,3,4,5,6,7,8,9,10,11,12,13,14,15,16,17},minor ytick={-1,0,1,2,3,4,5,6,7,8,9,10,11,12,13,14,15,16,17},grid=both,minor grid style={fmmuted!15},major grid style={fmmuted!15},tick label style={font=\tiny}]\fill[fmR,opacity=0.22] (axis cs:10.0000,0.0000)--(axis cs:4.0000,6.0000)--(axis cs:0.0000,8.0000)--(axis cs:0.0000,0.0000)--cycle;\draw[fmA,line width=1.5pt] (axis cs:-2.0000,9.0000)--(axis cs:18.0000,-1.0000);\draw[fmB,line width=1.5pt] (axis cs:-2.0000,12.0000)--(axis cs:18.0000,-8.0000);\draw[fmV,line width=1.5pt] (axis cs:0.0000,-2.0000)--(axis cs:0.0000,18.0000);\draw[fmK,line width=1.5pt] (axis cs:-2.0000,0.0000)--(axis cs:18.0000,0.0000);\end{axis}\end{tikzpicture}\end{center}

\fmm{F(0,0)=0;\ \ F(10,0)=300;\ \ F(4,6)=120+240=\mathbf{360};\ \ F(0,8)=320}

El beneficio máximo es de \textbf{360 €} y se obtiene haciendo \textbf{4 coches y 6 trenes} al día.
\end{enumerate}
\end{fmex}

\Needspace{5\baselineskip}
\begin{fmex}{Ejemplo resuelto: un problema de mínimo}
\begin{enumerate}[label=\textbf{\arabic*.},leftmargin=*]
\item Un ganadero mezcla dos piensos, A y B. Cada kg de A tiene 1 unidad de proteínas y 2 de vitaminas, y cada kg de B, 2 de proteínas y 1 de vitaminas. Cada día necesita como mínimo 8 unidades de proteínas y 10 de vitaminas. El pienso A cuesta 3 €/kg y el B, 2 €/kg. ¿Qué mezcla es la más barata?

\fmm{x=} kg de A, \fmm{y=} kg de B. \fmm{\ \begin{cases}x+2y\ge8\\ 2x+y\ge10\\ x\ge0,\ y\ge0\end{cases}\qquad C(x,y)=3x+2y} (mínimo)

La región no es acotada (hacia arriba y a la derecha). Vértices: \fmm{(0,10)}, \fmm{(8,0)} y el corte de las dos rectas, \fmm{(4,2)}.

\begin{center}\begin{tikzpicture}\begin{axis}[fmaxis,axis equal image,width=4.6cm,height=4.6cm,xmin=-1,xmax=12,ymin=-1,ymax=12,xtick={2,4,6,8,10,12},ytick={2,4,6,8,10,12},minor xtick={-1,0,1,2,3,4,5,6,7,8,9,10,11,12},minor ytick={-1,0,1,2,3,4,5,6,7,8,9,10,11,12},grid=both,minor grid style={fmmuted!15},major grid style={fmmuted!15},tick label style={font=\tiny}]\fill[fmR,opacity=0.22] (axis cs:12.0000,0.0000)--(axis cs:12.0000,12.0000)--(axis cs:0.0000,12.0000)--(axis cs:0.0000,10.0000)--(axis cs:4.0000,2.0000)--(axis cs:8.0000,0.0000)--cycle;\draw[fmA,line width=1.5pt] (axis cs:-2.0000,5.0000)--(axis cs:13.0000,-2.5000);\draw[fmB,line width=1.5pt] (axis cs:-2.0000,14.0000)--(axis cs:13.0000,-16.0000);\draw[fmV,line width=1.5pt] (axis cs:0.0000,-2.0000)--(axis cs:0.0000,13.0000);\draw[fmK,line width=1.5pt] (axis cs:-2.0000,0.0000)--(axis cs:13.0000,0.0000);\end{axis}\end{tikzpicture}\end{center}

\fmm{C(0,10)=20;\ \ C(4,2)=12+4=\mathbf{16};\ \ C(8,0)=24}

Como la región se extiende hacia donde \fmm{C} crece, el mínimo existe: la mezcla más barata es de \textbf{4 kg de A y 2 kg de B}, que cuesta \textbf{16 €}.
\end{enumerate}
\end{fmex}

\Needspace{5\baselineskip}
\begin{fmprop}{Rectas de nivel y casos especiales}
\begin{itemize}
\item Las rectas \fmm{ax+by=k} (todas paralelas) son las \textbf{rectas de nivel}: en todos sus puntos \fmm{F} vale \fmm{k}. Si desplazamos la recta en el sentido en que crece \fmm{F}, la última que toca la región marca el máximo.
\item \textbf{Infinitas soluciones:} si la recta de nivel es paralela a un lado, el óptimo se alcanza en todo el lado. Por ejemplo, maximizar \fmm{F=x+y} en la región de los juguetes: \fmm{F(10,0)=F(4,6)=10}.
\item \textbf{Sin óptimo:} en una región no acotada, \fmm{F} puede crecer tanto como queramos. Por ejemplo, maximizar \fmm{F=x+y} en la región de los piensos.
\end{itemize}
\end{fmprop}

  
\end{document}
